% ECONOMICS_PROBLEM: {"id":"C63-35","title":"Scalable heterogeneous-agent macroeconomics","jel_code":"C63","source_id":"OP-0016"}
% FIRST_POSTED: 2026-10-10
% ATTEMPT_STATUS: solved
% ENTRY_KIND: research_attempt
\documentclass[11pt,a4paper]{article}
\usepackage[left=20mm,right=20mm,top=18mm,bottom=18mm]{geometry}
\usepackage[T1]{fontenc}
\usepackage{lmodern,microtype}
\usepackage{amsmath,amssymb,amsthm,mathtools}
\usepackage{xurl,hyperref}
\hypersetup{hidelinks,hypertexnames=false,pdftitle={C63-35: Certified global nonlinear heterogeneous-agent computation}}
\setlength{\parindent}{0pt}
\setlength{\parskip}{3pt}
\setlength{\emergencystretch}{2em}
\allowdisplaybreaks
\raggedbottom
\widowpenalty=10000
\clubpenalty=10000
\displaywidowpenalty=10000
\newtheorem{theorem}{Theorem}
\newtheorem{lemma}[theorem]{Lemma}
\newtheorem{proposition}[theorem]{Proposition}
\newtheorem{corollary}[theorem]{Corollary}
\newcommand{\R}{\mathbb R}
\newcommand{\E}{\mathbb E}
\newcommand{\N}{\mathbb N}
\title{C63-35: Certified global nonlinear\\heterogeneous-agent computation}
\author{Ji Ho Bae}
\date{10 October 2026}
\begin{document}
\maketitle
\begin{abstract}
For an explicit nonempty class of heterogeneous-agent economies, we construct
and verify a competitive equilibrium on the full probability-measure state
domain, including transitions between binding asset constraints. Endogenous
prices, profits and distribution evolution enter a globally contractive Euler
update whose constants are derived from the primitives. Exact piecewise-affine
inversion gives a finite global solver with polynomial real-arithmetic cost in
inverse accuracy for the fixed shock-state count. A separately certified finite
precision implementation has an explicit quasipolynomial bit bound. Economic
residuals control distance to an actual equilibrium. A concrete identified
Gaussian experiment propagates numerical error through estimation, likelihoods
and equilibrium counterfactual confidence intervals. Reproducible remote
benchmarks measure policy and complete value queries, including preprocessing
and actual exact-rational bit lengths.
\end{abstract}

\section*{Scope and process}
\textbf{Model:} GPT Codex. \textbf{Contributor:} Ji Ho Bae.
The exact serving revision was not exposed. The construction and hand proof
were produced by parallel derivation and independent adversarial review.
The catalogue asks for a declared restricted primitive class $\mathcal H$,
a global residual guarantee, stability, estimation propagation and a controlled
benchmark. We specify every member of $\mathcal H$ below and prove these
requirements on its full state domain. The announced industry-price selection
is explicit; other competitive boundary-price selections may exist.

Heterogeneous-agent saving and aggregate-risk computation have established
benchmarks~\cite{Aiyagari,KS,HANK}. The sequence-space approach supplies efficient
solution and estimation tools, including nonlinear perfect-foresight
transitions~\cite{ABRS}. Our storage economy, global certificate and local
Taylor comparison are the constructions specified here. The comparison is
not an implementation of those authors' algorithms. All mathematical steps
needed for the result are proved below; citations acknowledge this background.
This is an AI-assisted hand proof, with no proof-assistant or historical-priority
claim; community expert review remains separate.

The complete original generated output is preserved at the immutable
\href{https://github.com/jbaelaw/certified-heterogeneous-agent-computation/blob/629e26c2cc4fce0c729a0545c296e537fcd38491/paper/Proof.tex}{manuscript TeX}
and \href{https://github.com/jbaelaw/certified-heterogeneous-agent-computation/blob/629e26c2cc4fce0c729a0545c296e537fcd38491/paper/Proof.pdf}{manuscript PDF}.

\paragraph{Notation.}
The claim price is $q=\beta Q$, the technological marginal-cost intercept is
$q_0=4/5$, and $\rho$ is the coupled policy-update modulus, bounded uniformly
by $\bar\rho=0.021$. Input representations, arithmetic and bit costs
are distinguished in Section~\ref{sec:finite-solver}.

% Body component: concrete global heterogeneous-agent economy and its proof.
% Other agents own solver/inference/audit components. This file alone is not
% a completion claim for the whole C63-35 statement.
\section{A concrete primitive class on the full distribution domain}
\label{sec:economy}

Let $E=\{L,H\}$, assets lie in $A=[0,1]$, and $Z$ is a fixed two-element
aggregate-shock set with a declared Markov matrix $P$ whose entries lie
in $[0.1,0.9]$. Every agent
observes its own asset and current idiosyncratic state and the public
aggregate state $(z,\mu)$. Put
\[
 \mathcal D=A\times E\times Z\times\mathcal P(E\times A),\qquad
 d((e,a),(f,b))=|a-b|+\mathbf1\{e\ne f\},
 \tag{EC1}
\]
and equip population laws with the corresponding $1$-Wasserstein distance
$W_1$. No restriction to stationary or simulated population laws is made.
Next period's idiosyncratic state is a fair independent draw from $E$,
independent of assets and aggregate shocks. Population laws evolve by this
kernel exactly, rather than by a finite-sample approximation. One concrete
population realization uses initial law $\mu$ times a countable product of
uniform coordinates, splitting each state cohort equally by each successive
coordinate; future coordinates are not in households' information sets.

The primitive class has
\[
\begin{gathered}
 0.005\leq\beta\leq0.01,\quad 0.0001\leq\kappa\leq0.001,\quad q_0=4/5,\\
 y_L(z)=1+\delta_z,\quad y_H(z)=3+\delta_z,\quad |\delta_z|\leq1/10.
\end{gathered}
 \tag{EC2}
\]
Household utility is $u(c)=5c-c^2/2$. An asset is a one-period claim to one
unit of next period's consumption good. A household begins with $a$ mature
claims, backed by installed output; the initial stock is an endowment.
There is no other asset, borrowing or storage. If the new-claim price is
$q=\beta Q$ and the equal per-household current dividend is $\pi$, the
budget and choice are
\[
 c=y_e(z)+a+\pi-\beta Qb,\qquad b\in[0,1].
 \tag{EC3}
\]
An atomless household takes aggregate states and prices as given.

A competitive investment industry sells $K$ new claims for $qK$, uses
$\beta C(K)$ current goods, and produces $K$ goods next period, where
$C(K)=q_0K+\kappa K^2/2$. Households own equal shares of its profits.
We use the explicit marginal-cost price selection
\[
 K=\int g(a,e,z,\mu)\,d\mu(e,a),\qquad
 Q=q_0+\kappa K,\qquad
 \pi=qK-\beta C(K)=\frac{\beta\kappa K^2}{2}.
 \tag{EC4}
\]
It is a competitive selection: $K$ maximizes
$qk-\beta C(k)$ over $k\geq0$. At $K=0$ a competitive firm alone also
permits prices below marginal cost; equality in \textup{(EC4)} is the
announced selection there, not a consequence claimed from an interior
first-order condition.

For every household policy $g$ define $K_g,Q_g,\pi_g$ by
\textup{(EC4)} and
\[
 \Phi_g(z,\mu)=\mathcal L\bigl(e',g(a,e,z,\mu)\bigr),
 \qquad (e,a)\sim\mu,\quad e'\sim\operatorname{Unif}(E)
 \text{ independently}.
 \tag{EC5}
\]
The next aggregate state is $(z',\Phi_g(z,\mu))$ with $z'\sim P(z,\cdot)$.
This is the endogenous full population transition. Integrating
\textup{(EC3)} verifies goods clearing exactly:
\[
 \int c\,d\mu+\beta C(K_g)=\int y_e(z)\,d\mu+\int a\,d\mu.
 \tag{EC6}
\]
Claims clearing is $K_g=\int g\,d\mu$, and next period's mature output
is $K_g=\int a'\,d\Phi_g$. Thus profits are counted once as household
income and once as the industry's net revenue, not as extra resources.

Let $\mathcal S$ consist of policies $g:\mathcal D\to[0,1]$ satisfying
\[
 |g(a,e,z,\mu)-g(b,e,z,\mu)|\leq|a-b|,
 \quad |g(a,e,z,\mu)-g(a,e,z,\nu)|\leq W_1(\mu,\nu).
 \tag{EC7}
\]
They are continuous in assets and population laws. Different idiosyncratic
labels have distance one, and boundedness makes the policy $1$-Lipschitz
in $(e,a)$ as well. With the uniform norm, $\mathcal S$ is nonempty,
closed and complete. The domain includes every probability law in
\textup{(EC1)}. Throughout put
\[
 Q_{\max}=0.801,\quad \pi_{\max}=0.000005,
 \quad \lambda=1-\beta Q_{\max}\geq0.99199.
 \tag{EC8}
\]
Every feasible household choice under the induced prices has
\[
 0.89199\leq c\leq4.100005<5,
 \qquad 0<u'(c)\leq4.1+\beta Q_{\max},\qquad |u(c)|\leq12.5.
 \tag{EC9}
\]
Consequently every adapted asset plan is consumption-feasible and every
discounted objective is bounded.

\section{A global Euler update for the declared selection}
\label{sec:euler-update}
For $g\in\mathcal S$ and current state $(a,e,z,\mu)$ put
$\mu'=\Phi_g(z,\mu)$ and define
\begin{align*}
 A_g(b;z,\mu)
 &=\sum_{z'}\frac{P_{zz'}}2\sum_{e'\in E}
 \bigl[5-y_{e'}(z')-\pi_g(z',\mu')
       +\beta Q_g(z',\mu')g(b,e',z',\mu')\bigr],\\
 D_g(b;a,e,z,\mu)
 &=A_g(b;z,\mu)-Q_g(z,\mu)
       [5-y_e(z)-a-\pi_g(z,\mu)]
       -(1+\beta Q_g(z,\mu)^2)b.
 \tag{EC10}
\end{align*}
The update $F_g(a,e,z,\mu)$ is the unique $b\in[0,1]$ satisfying
\[
 b=\operatorname{clip}_{[0,1]}
 \frac{A_g(b;z,\mu)-Q_g(z,\mu)[5-y_e(z)-a-\pi_g(z,\mu)]}
      {1+\beta Q_g(z,\mu)^2}.
 \tag{EC11}
\]
It is a future-policy Euler update, not the household best-response map
for an arbitrary candidate $g$. Household optimality is proved below only
at a fixed point.

\begin{lemma}[Well-defined update and its root bound]\label{lem:euler-root}
The inner map in \textup{(EC11)} contracts in $b$ with constant at most
$\beta Q_{\max}<1$. Equivalently $b=F_g$ obeys
\[
 D_g(b)\begin{cases}\leq0,&b=0,\\=0,&0<b<1,\\\geq0,&b=1.\end{cases}
 \tag{EC12}
\]
For $b'>b$, $D_g(b')-D_g(b)\leq-\lambda(b'-b)$.
If two such decreasing functions differ uniformly by at most $\eta$,
their constrained roots differ by at most $\eta/\lambda$.
\end{lemma}
\begin{proof}
The $b$-Lipschitz bound for $A_g$ is $\beta Q_{\max}$ by
\textup{(EC7)}; its other terms do not depend on the individual asset
choice, since an atomless household cannot change $\mu'$. The clipping
map is $1$-Lipschitz. This proves inner contraction and strict decrease
of $D_g$ with the stated modulus. Clipping gives exactly
\textup{(EC12)}. To prove the root bound, suppose roots $b_1>b_2$.
Then the first function is nonnegative at $b_1$, hence at least
$\lambda(b_1-b_2)$ at $b_2$, while the second is nonpositive at $b_2$.
Their uniform difference bounds $\lambda(b_1-b_2)$ by $\eta$.
The reverse ordering is identical.
\end{proof}

\begin{theorem}[Global self-mapping and contraction]\label{thm:economy-contraction}
$F$ maps $\mathcal S$ into itself. More precisely,
\begin{align*}
 \operatorname{Lip}_a(F_g)&\leq Q_{\max}/\lambda<0.808,\\
 \operatorname{Lip}_\mu(F_g)&\leq
 \frac{8.2\kappa+\beta[(6Q_{\max}+8)\kappa+2Q_{\max}]}\lambda<0.025,\\
 \|F_g-F_h\|_\infty&\leq \rho\|g-h\|_\infty,\qquad
 \rho:=\frac{4.1\kappa+\beta[(3Q_{\max}+6)\kappa+2Q_{\max}]}\lambda<0.021.
 \tag{EC13}
\end{align*}
It has one fixed point $g^*\in\mathcal S$ on the entire domain
\textup{(EC1)}.
\end{theorem}
\begin{proof}
For $d=\|g-h\|_\infty$ at the same $(z,\mu)$,
\[
 |K_g-K_h|\leq d,\quad |Q_g-Q_h|\leq\kappa d,
 \quad |\pi_g-\pi_h|\leq\beta\kappa d,
 \quad W_1(\Phi_g,\Phi_h)\leq d.
 \tag{EC14}
\]
For one policy at population laws distance $w$, integration of its
$1$-Lipschitz state function and \textup{(EC7)} gives
\[
 |K_g(z,\mu)-K_g(z,\nu)|\leq2w,
 \qquad W_1(\Phi_g(z,\mu),\Phi_g(z,\nu))\leq2w.
 \tag{EC15}
\]
In \textup{(EC14)}, the next-state $K$ difference is therefore at most
$3d$; the next-policy difference at fixed $b,e',z'$ is at most $2d$.
Next dividends differ by at most $3\beta\kappa d$, and next prices by
at most $3\kappa d$. Thus the $A$-difference is at most
$\beta(6\kappa+2Q_{\max})d$. The current product $Q[5-y-a-\pi]$
differs by at most $(4.1\kappa+\beta Q_{\max}\kappa)d$, because
$0<5-y-a-\pi\leq4.1$. The coefficient $\beta Q^2b$ contributes at most
$2\beta Q_{\max}\kappa d$. Adding proves the numerator of $\rho$.
Apply Lemma~\ref{lem:euler-root}.

For population variation, next-state $K$ differs by at most $4w$ and
next-policy values by at most $2w$. The corresponding $A$ bound is
$\beta(8\kappa+2Q_{\max})w$. The current product and quadratic
coefficient contribute respectively
$(8.2\kappa+2\beta Q_{\max}\kappa)w$ and
$4\beta Q_{\max}\kappa w$. Their sum gives the asserted population
modulus. Asset variation contributes only $Q_g|a-b|$ to $D_g$.
All three inequalities are global. Substituting the endpoint bounds in
\textup{(EC2)} and \textup{(EC8)} gives the displayed numerical upper
bounds. Clipping gives range $[0,1]$, so $F_g\in\mathcal S$.

Starting from any $g_0\in\mathcal S$, consecutive iterates have distances
at most $\rho^k\|F_{g_0}-g_0\|_\infty$. Their geometric tail proves a
uniform Cauchy limit in the complete set $\mathcal S$. Continuity of $F$
gives a fixed point. Two fixed points have distance at most $\rho$ times
that distance, hence coincide. No market-clearing contraction has been
assumed; its constants have just been derived from the stated primitives.
\end{proof}

\section{Household verification, markets and binding constraints}
\begin{theorem}[Actual competitive equilibrium]\label{thm:actual-equilibrium}
The policy $g^*$, prices and dividends \textup{(EC4)}, and transition
\textup{(EC5)} form a recursive competitive equilibrium under the
marginal-cost selection. Household optimization holds against all adapted,
possibly randomized asset plans, at every state in $\mathcal D$.
Conversely every such selected recursive equilibrium with policy in
$\mathcal S$ is this fixed point. No uniqueness outside $\mathcal S$ or
across other competitive boundary-price selections is asserted.
\end{theorem}
\begin{proof}
Fix any initial state and let $(a_t^*,e_t,z_t,\mu_t)$ follow $g^*$.
An individual deviation leaves $(z_t,\mu_t)$ and its prices unchanged.
Couple the same idiosyncratic and aggregate shocks, and write the deviating
assets as $a_t$, with $a_0=a_0^*$. All its choices in $[0,1]$ are feasible
by \textup{(EC9)}. Put $\Delta a_t=a_t-a_t^*$ and
$m_t=u'(c_t^*)$. Concavity gives
\[
 u(c_t)-u(c_t^*)\leq m_t[\Delta a_t-\beta Q_t\Delta a_{t+1}].
 \tag{EC16}
\]
Summing its discounted expectation through date $N-1$ yields
\[
 \sum_{t=1}^{N-1}\beta^t
 \E[(\E[m_t\mid\mathcal F_{t-1}]-Q_{t-1}m_{t-1})\Delta a_t]
 -\beta^N\E[Q_{N-1}m_{N-1}\Delta a_N].
 \tag{EC17}
\]
Here $\mathcal F_{t-1}$ includes the deviator's private randomization and
its date-$t-1$ choice; $\Delta a_t$ is measurable there. The conditional
marginal-utility difference is precisely $D_{g^*}$ evaluated at its
baseline choice $a_t^*$. By \textup{(EC12)}, its product with
$a_t-a_t^*$ is nonpositive: the difference is nonpositive at a lower
bound, nonnegative at an upper bound, and zero in the interior.
The final term in absolute value is at most
$\beta^NQ_{\max}(4.1+\beta Q_{\max})$, tending to zero.
Bounded discounted utility allows $N\to\infty$, proving optimality
against every adapted deviation. This does not differentiate a value
function or assume an envelope theorem across binding constraints.

Necessity of the Euler inequalities uses a two-period deviation.
Perturb the current choice by a small feasible amount, then at the next
date choose the baseline next asset irrespective of the perturbed asset.
Consumption remains positive by \textup{(EC9)}, and from the following
date onward the original asset path is restored. The directional
utility derivative is $\beta D_g$ times that perturbation. Optimality
therefore implies exactly \textup{(EC12)}, including its boundary signs.
This proves the converse through uniqueness of the constrained root and
Theorem~\ref{thm:economy-contraction}. Industry optimality and both goods
and claims clearing were verified in \textup{(EC4)--(EC6)}.
\end{proof}

Define the actual policy value
\[
 V_g(a,e,z,\mu)=\E\sum_{t\geq0}\beta^t u(c_t^g),
 \qquad \|V_g\|_\infty\leq12.5/(1-\beta).
 \tag{EC18}
\]
At $g=g^*$ this is the household optimal value and obeys the full Bellman
maximization equation: decomposition of \textup{(EC18)} gives equality
at $g^*$, and any current choice followed by $g^*$ is among the deviations
just bounded. Finite-horizon policy values are continuous, and their
uniformly convergent discounted limit is continuous. Thus the value and
Euler descriptions concern the same fully specified optimization problem.

\begin{proposition}[Both boundaries and an interior occur globally]\label{prop:actual-binding}
For every admissible primitive and aggregate state $(z,\mu)$,
\[
 g^*(0,L,z,\mu)=0,\qquad
 g^*(a,H,z,\mu)=1\ (0\leq a\leq1),\qquad
 0.2<g^*(0.75,L,z,\mu)<0.6.
 \tag{EC19}
\]
\end{proposition}
\begin{proof}
The same signs hold for every $F_g$, $g\in\mathcal S$. The average
$5-y_{e'}(z')$ is between $2.9$ and $3.1$. Consequently
$2.9-\pi_{\max}\leq A_g(b)\leq3.1+\beta Q_{\max}$.
For low income and zero assets,
\[
 D_g(0)\leq3.1+0.00801-0.8(3.9-0.000005)=-0.011986<0.
\]
For high income and any assets,
\[
 D_g(1)\geq2.899995-0.801(2.1)-(1+0.01\cdot0.801^2)
 =0.21147899>0.
\]
For low income and assets $0.75$,
$D_g(0)\geq2.899995-0.801(3.35)=0.216645$ and
$D_g(1)\leq3.10801-0.8(3.15-0.000005)-1=-0.411986$.
The absolute $b$-Lipschitz constant of $D_g$ is at most
$1+\beta Q_{\max}^2+\beta Q_{\max}\leq1.01442601$.
Thus $D_g(0.2)>0$ and $D_g(0.6)<0$, placing its root strictly between
them. Pass to the fixed point. Continuity in assets connects genuinely
binding and interior regimes throughout the full population domain.
\end{proof}

\section{Useful uniform output and primitive stability bounds}
\begin{proposition}\label{prop:economy-output-stability}
For $g,h\in\mathcal S$, $d=\|g-h\|_\infty$, at the same state
\[
 \|K_g-K_h\|_\infty\leq d,\quad
 \|Q_g-Q_h\|_\infty\leq\kappa d,\quad
 \|\pi_g-\pi_h\|_\infty\leq\beta\kappa d,\quad
 \sup_{z,\mu}W_1(\Phi_g,\Phi_h)\leq d.
 \tag{EC20}
\]
The policy values in \textup{(EC18)} obey
\[
\begin{gathered}
 \|V_g-V_h\|_\infty\leq L_Vd<0.08d,\\
 L_V=(4.1+\beta Q_{\max})
 \left[(1+2\beta Q_{\max}+4\beta\kappa)
 \left(\frac1{1-2\beta}-\frac1{1-\beta}\right)
 +\frac{\beta(Q_{\max}+2\kappa)}{1-\beta}\right].
\end{gathered}
 \tag{EC21}
\]
For fixed $\beta,P$ the equilibrium policy has the primitive sensitivities
\[
 \|g^*_{\kappa,\delta}-g^*_{\tilde\kappa,\delta}\|_\infty
\leq\frac{4.1+\beta(2.5Q_{\max}+1.5)}{\lambda(1-\bar \rho)}
 |\kappa-\tilde\kappa|<5|\kappa-\tilde\kappa|,
 \tag{EC22}
\]
\[
 \|g^*_{\kappa,\delta}-g^*_{\kappa,\tilde\delta}\|_\infty
\leq\frac{1+Q_{\max}}{\lambda(1-\bar \rho)}
 \|\delta-\tilde\delta\|_\infty<2\|\delta-\tilde\delta\|_\infty.
 \tag{EC23}
\]
Here $\bar \rho=0.021$ is the common uniform contraction bound when
comparing different primitives.
\end{proposition}
\begin{proof}
The first inequalities are \textup{(EC14)}. For values, couple the same
$e,z$ shocks from a common initial state. If $r_t$ is population $W_1$
distance and $b_t$ individual asset distance, then
$r_{t+1}\leq2r_t+d$ and $b_{t+1}\leq b_t+r_t+d$.
Both start at zero, hence both are at most $(2^t-1)d$.
At the coupled states current prices and dividends differ by at most
$\kappa(2r_t+d)$ and $\beta\kappa(2r_t+d)$.
Using the budgets, the consumption difference is at most
\[
 (1+\beta Q_{\max})b_t+
 \beta(Q_{\max}+4\kappa)r_t+\beta(Q_{\max}+2\kappa)d.
\]
The bound on $u'$ in \textup{(EC9)} and the geometric sums
$\sum\beta^t(2^t-1)=(1-2\beta)^{-1}-(1-\beta)^{-1}$ and
$\sum\beta^t=(1-\beta)^{-1}$ prove \textup{(EC21)}.

For one fixed candidate policy, changing $\kappa$ does not change
$\Phi_g$. It changes $Q$ by at most $|\Delta\kappa|$ and $\pi$ by
at most $\beta|\Delta\kappa|/2$. In \textup{(EC10)} the current product,
quadratic coefficient and future bracket together change by at most
$[4.1+\beta(2.5Q_{\max}+1.5)]|\Delta\kappa|$.
Changing income offsets changes only current and expected future income,
so its effect is at most $(1+Q_{\max})\|\Delta\delta\|_\infty$.
Apply the root bound, then write the difference of the two fixed points
as a same-primitive contraction difference plus this primitive difference.
Moving its $\rho$ term to the left gives \textup{(EC22)--(EC23)}.
\end{proof}

This is a nonlinear heterogeneous-agent equilibrium with endogenous prices,
profits and a full distribution transition. For example set every
$\delta_z=0$. The reference policy
$g_{\rm aux}(a,L)=\max\{0,0.8a-0.2\}$, $g_{\rm aux}(a,H)=1$ has the formal
zero-$\beta$, zero-$\kappa$ Euler root. For the positive primitives in
\textup{(EC2)}, direct comparison of \textup{(EC10)} with that scalar
root gives
\[
 \|g^*-g_{\rm aux}\|_\infty
 \leq\frac{4.1\kappa+\beta Q_{\max}\kappa/2+
                  \beta\kappa/2+\beta Q_{\max}+\beta Q_{\max}^2}
                 {\lambda}<0.019.
 \tag{EC24}
\]
Thus the two initial populations concentrated on low income, respectively
at asset $1/2$ and equally at assets $0,1$, have the same asset mean and
type marginal but investment demands differing by at least
$0.1-2(0.019)>0.06$. Their selected prices therefore differ because
$\kappa>0$. This comparison uses the limiting formula only as an explicit
auxiliary root; zero discounting is not included in the primitive class.
The full population state cannot silently be replaced by its mean.

% Root-owned body fragment. Does not alone complete the benchmark/inference task.
\section{An economic residual and a linear error bound}
\label{sec:residual-stability}

Here is a residual norm for an explicitly specified candidate format.
Let $g\in\mathcal S$, let $\widehat K(z,\mu)\in[0,1]$, let
$\widehat\Phi(z,\mu)$ be a probability law, and let $\widehat V$ be
bounded. All these maps are measurable. Set
$\widehat Q=q_0+\kappa\widehat K$ and
$\widehat\pi=\beta\kappa\widehat K^2/2$, and define consumption by
the household budget with these prices. This format enforces the
announced price selection and the industry's optimizing supply choice;
it does not assume that capital demand or population evolution is correct.
Use the maximum uniform norm for scalar functions and the uniform $W_1$
distance for population transitions. The full output tuple comprises
$(g,K,Q,\pi,\Phi,V)$ in this norm.

Put $K_g=\int g\,d\mu$ and define
\[
 r_K=\|\widehat K-K_g\|_\infty,\qquad
 r_\Phi=\sup_{z,\mu}W_1(\widehat\Phi(z,\mu),\Phi_g(z,\mu)).
 \tag{R1}
\]
Let $\widehat D_g(b)$ be the Euler numerator \textup{(EC10)}, using
$\widehat Q,\widehat\pi,\widehat\Phi$ in place of their induced
counterparts. The saving-complementarity residual is
\[
 r_E=\sup_{\mathcal D}
 \left|g-\operatorname{clip}_{[0,1]}(g+\widehat D_g(g))\right|.
 \tag{R2}
\]
It vanishes for the correct interior equality and for each satisfied
boundary inequality, including strict slack. Let $\widehat T_g$ be the
one-period policy-value operator at the returned prices and transition,
and let $\widehat B$ denote the pointwise supremum over all $b\in[0,1]$; no attainment is needed for a merely measurable candidate value. Include
\[
 r_V=\|\widehat V-\widehat T_g\widehat V\|_\infty,
 \qquad
 r_{\rm opt}=\|[\widehat B\widehat V-\widehat V]_+\|_\infty.
 \tag{R3}
\]
The residual $\|\mathcal R\|$ is the maximum of these quantities,
the absolute goods-market error, and any budget or feasibility violations.
Budget and saving feasibility hold exactly in this format. Firm optimality
holds exactly by the marginal-cost choice. Thus this is an economic
residual, with zero contribution from satisfied slack inequalities.

\begin{proposition}[Global residual-to-equilibrium control]
\label{prop:full-residual-stability}
For the stated format and every primitive in \textup{(EC2)},
\[
 \operatorname{dist}(\widehat f,\mathcal E)
 \leq\|\widehat f-f^*\|\leq3\|\mathcal R(\widehat f)\|.
 \tag{R4}
\]
Here $f^*$ is the constructed selected equilibrium and $\mathcal E$
may contain other competitive equilibria. No uniqueness of that larger
set is required. The bound is uniform over the full distribution domain.
\end{proposition}
\begin{proof}
Write $r=\|\mathcal R\|$ and $\bar\rho=0.021$. At a common state,
\[
 |\widehat Q-Q_g|\leq\kappa r,\qquad
 |\widehat\pi-\pi_g|\leq\beta\kappa r.
\]
At the two next population laws, the corresponding differences are
at most $3\kappa r$ and $3\beta\kappa r$, because
$K_g$ has distribution Lipschitz constant at most two. The future
policy difference at a fixed individual asset is at most $r$.
Expanding \textup{(EC10)} therefore gives
\[
 \|\widehat D_g-D_g\|_\infty
 \leq\left[4.1\kappa+
 \beta\{(3Q_{\max}+6)\kappa+Q_{\max}\}\right]r<0.013r.
 \tag{R5}
\]

For any decreasing Euler numerator with decrease modulus $\lambda$,
let $b^*$ be its constrained root. At $b>b^*$ its numerator is at
most $-\lambda(b-b^*)$. Consequently
$|b-\operatorname{clip}(b+D(b))|
 =\min\{b,-D(b)\}\geq\lambda(b-b^*)$.
At $b<b^*$ the symmetric argument applies, including an upper-bound
root. Since $\lambda\leq1$, the clipping endpoints do not weaken this
bound. The clipping map is also $1$-Lipschitz. By \textup{(R2),(R5)},
\[
 \|g-F_g\|_\infty\leq\frac{1.013r}{\lambda},\qquad
 d:=\|g-g^*\|_\infty
 \leq\frac{1.013r}{\lambda(1-\bar\rho)}<1.05r.
 \tag{R6}
\]
Capital and population-transition errors relative to equilibrium are
at most $r+d$; the price and dividend errors are smaller multiples of
that bound.

Compare $\widehat T_g$ with the exact tuple's policy operator $T_g$.
Current utility changes by at most $10\beta\kappa r$, and the
continuation of the exact policy value $V_g$ changes by at most
$6\beta r$. The Lipschitz bound $6$ follows by coupling the same
shocks: consumption has Lipschitz constants at most $1.009$ in assets
and $0.0081$ in population, the joint next-state distance grows by at
most three, and $5(1.009)/(1-3\beta)<6$.
The policy operator contracts bounded values by $\beta$. Thus
\[
 \|\widehat V-V_g\|_\infty
 \leq\frac{1+10\beta\kappa+6\beta}{1-\beta}\,r<1.072r.
\]
Equation \textup{(EC21)} and \textup{(R6)} then give
$\|\widehat V-V^*\|_\infty<1.072r+0.08(1.05r)<1.16r$.
Every tuple component is therefore within $3r$ of $f^*$. The latter
is a competitive equilibrium, so distance to $\mathcal E$ is no larger.
The optional Bellman-optimality and resource blocks only increase $r$.
\end{proof}

In particular, zero residual in this format is equivalent to an actual
selected equilibrium, with the correct policy value. This does not
require an Euler equality at a binding household constraint, or assert
that the selected industry price is the only competitive price at zero
production.

\section{A finite solver on the entire distribution state space}
\label{sec:finite-solver}

We now construct the evaluation procedure, rather than assume a certified
equilibrium-evaluation oracle. Use the economy and operator from the
preceding construction. Put
\[
 Q_{\max}=0.801,\quad m_0=0.009,\quad \bar\rho=0.021,\quad
 L=|Z|,\qquad
 \mathcal D=[0,1]\times E\times Z\times\mathcal P(E\times[0,1]).
 \tag{S1}
\]
The metric on $E\times[0,1]$ is
$\mathbf1_{\{e\ne e'\}}+|a-a'|$ and its Wasserstein distance is $W_1$.
The policy class $\mathcal S$ consists of functions valued in $[0,1]$
with asset and distribution Lipschitz constants at most one, separately
for each finite $(e,z)$. Its distance is the uniform norm on
\emph{all} of $\mathcal D$.

\subsection{The finite input representation}
There is no numerical algorithm whose input is an unencoded arbitrary
probability measure. A numerical query specifies $(a,e,z)$ and either a
finite weighted empirical distribution, or the finite barycentric
features of $\mu$ used below. In the latter representation, each feature
is supplied as an exact real-arithmetic input, or by an oracle returning
a certified interval of any requested rational width; each query and its
precision cost is charged. Such an oracle represents input data only,
not a household solution or an equilibrium evaluation. The global output
is a finite program defining a function of these features for every
$\mu\in\mathcal P(E\times[0,1])$. Its uniform certificate holds
mathematically for every such $\mu$. Numeric evaluation is claimed only
under the declared input representation.

\subsection{A nonexpansive distribution encoding}
For integer $M\geq1$, let $a_j=j/M$ and
$\lambda_j(a)=(1-M|a-a_j|)_+$, $0\leq j\leq M$. Define
\[
 Q_M\mu(\{(e,a_j)\})=\int_{\{e\}\times[0,1]}\lambda_j(a)\,\mu(de,da).
 \tag{S2}
\]
The hat functions sum to one. Equivalently a point in
$[a_j,a_{j+1}]$ is rounded to its two endpoints with probabilities
preserving its mean. Thus $Q_M$ is a Markov map, fixes grid-supported
measures, and uses $N=2(M+1)$ nonnegative weights summing to one.
It satisfies
\[
 W_1(Q_M\mu,Q_M\nu)\leq W_1(\mu,\nu),\qquad
 W_1(Q_M\mu,\mu)\leq\frac1{2M}.
 \tag{S3}
\]
Indeed the rounded laws of two asset points are stochastically ordered
in the same order as those points, and have their respective means.
Coupling them by common quantiles therefore gives expected asset
distance exactly $|a-b|$. Keep the discrete type coordinate unchanged
and apply this coupling to any coupling of the original measures.
For the second bound, a point a fraction $r$ across its grid interval
has expected displacement $2r(1-r)/M\leq1/(2M)$.
This encoding is continuous in $W_1$; hard cell assignment would not
have that property.

\subsection{A finite operator that remains in the certified domain}
For $g\in\mathcal S$ and a current state $\mu$, set
\[
 \nu=Q_M\mu,\qquad
 \nu'=Q_M\Phi_g(z,\nu),\quad
 Q=Q_g(z,\nu),\quad \pi=\pi_g(z,\nu).
 \tag{S4}
\]
All these quantities use finite sums over the $N$ grid states.
For a prospective saving choice $b$ define
\[
 \bar E_g(b)=
 \sum_{z'}P_{zz'}\frac12\sum_{e'}
 \left[5-y_{e'}(z')-\pi_g(z',\nu')
       +\beta Q_g(z',\nu')g(b,e',z',\nu')\right],
 \tag{S5}
\]
and the clipped map
\[
 \bar H_g(b)=
 \left[\frac{\bar E_g(b)-Q\{5-y_e(z)-a-\pi\}}
 {1+\beta Q^2}\right]_{[0,1]}.
 \tag{S6}
\]
Let $\bar F_Mg$ be its unique fixed point. The next subsection supplies
an exact finite piecewise-affine algorithm for all the iterates used
here. Every policy evaluation and distribution integral in that
algorithm is a finite computation; no root-evaluation oracle is assumed.

Here are explicit bounds for these maps. The direct dependence of
$\bar H$ on $b$ is at most $m_0$. For fixed $b$, its asset Lipschitz
constant is at most $Q_{\max}$; its distribution Lipschitz constant is
at most
\[
 C_\mu=
 8.2\kappa+\beta\{(6Q_{\max}+8)\kappa+2Q_{\max}\}
 \leq0.02434806.
 \tag{S7}
\]
For two policies its fixed-$b$ difference is at most
\[
 C_g\|g-h\|_\infty,\qquad
 C_g=4.1\kappa+
 \beta\{(3Q_{\max}+6)\kappa+2Q_{\max}\}
 \leq0.02020403.
 \tag{S8}
\]
These are the same elementary kernel and price estimates as in the
economic construction: $Q_M$ is nonexpansive, so it increases none of
them. One can verify the quotient step without assuming its numerator
lies in $[0,1]$: for $d,d'\geq1$,
\[
 |[n/d]_{[0,1]}-[n'/d']_{[0,1]}|
 \leq |n-n'|+|d-d'|.
 \tag{S9}
\]
To prove this, first compare the numerators with the denominator fixed,
and then compare denominators with the numerator fixed, splitting the
three clipped regions. In the interior the derivative in the denominator
has absolute value at most one; the boundary pieces preserve that bound.
The denominator contribution in \textup{(S7)} is
$4\beta Q_{\max}\kappa W_1(\mu,\tilde\mu)$ and in
\textup{(S8)} it is $2\beta Q_{\max}\kappa\|g-h\|_\infty$.

Iteration of $\bar H$ from zero and passage to its uniform limit
give asset Lipschitz constant at most
$Q_{\max}/(1-m_0)<0.809$ and distribution Lipschitz constant at most
$C_\mu/(1-m_0)<0.025$. Its values lie in $[0,1]$. Thus
\[
 \bar F_M:\mathcal S\longrightarrow\mathcal S,\qquad
 \|\bar F_Mg-\bar F_Mh\|_\infty
 \leq\frac{C_g}{1-m_0}\|g-h\|_\infty
 <\bar\rho\|g-h\|_\infty .
 \tag{S10}
\]
The finite-root construction below represents this limit exactly.
The approximate outer iterates therefore remain inside the domain
where the actual coupled equilibrium contraction was proved.

Let $F$ denote that exact equilibrium operator. Its distribution
Lipschitz constant is less than $0.025$. Projection of the current state
therefore changes $Fg$ by at most $0.025/(2M)$. Projection of the future
state changes the numerator of its clipped equation by at most
$\beta(4\kappa+Q_{\max})/(2M)$. The corresponding root changes by at most
that number divided by $1-m_0$. Hence
\[
 \|\bar F_Mg-Fg\|_\infty\leq\delta_M:=\frac{0.02}{M}.
 \tag{S11}
\]
No measure-state net or simulated-path coverage enters this bound.

\subsection{Exact roots by finite piecewise-affine inversion}
For a fixed grid distribution $\nu$ and aggregate state $z$, construct
both functions $a\mapsto g_k(a,e,z,\nu)$ together. They are continuous
piecewise affine. At level zero they are identically zero. Suppose this
has been proved and implemented at level $k-1$.

First construct the two level-$(k-1)$ functions at the current context
$(z,\nu)$. Evaluate them at the grid points to compute $K$, $Q$, $\pi$
and the rounded future distribution
$\nu'=Q_M\Phi_{g_{k-1}}(z,\nu)$. These quantities do not depend on the
individual current asset or type. Next construct the two
level-$(k-1)$ functions at each future context $(z',\nu')$.
The finite weighted sum \textup{(S5)} is piecewise affine in $b$,
with knots among the union of those $2L$ functions' knots.

For each current type $e$, the un-clipped Euler equation can be
written as
\[
 a=A_e(b):=
 \frac{(1+\beta Q^2)b-\bar E_{g_{k-1}}(b)
       +Q\{5-y_e(z)-\pi\}}{Q}.
 \tag{S11a}
\]
The function $A_e$ is continuous piecewise affine and strictly
increasing: every secant slope is at least
$(1-\beta Q_{\max})/Q_{\max}>0$. It can therefore be inverted on
its range one affine segment at a time. Set the answer to zero for
$a\leq A_e(0)$ and to one for $a\geq A_e(1)$, and restrict it to
$a\in[0,1]$. This is exactly $g_k=\bar F_Mg_{k-1}$, including both
binding regions. No root search or rounding of the root is used.
If every preceding policy has at most $B_{k-1}$ affine pieces,
the resulting policy has at most $2L B_{k-1}+2$ pieces. The safe
uniform bound
\[
 B_k=(2L+3)^{k+1}
 \tag{S11b}
\]
suffices at every context.

All computations are arithmetic, comparisons and sorting of finite
lists. Rational input primitives and grid-distribution weights give
rational knots, slopes and intercepts; exact rational arithmetic is
therefore an implementation. The division by $Q$ is safe because
$Q\geq0.8$, and each segment slope of $A_e$, which is a denominator
in the inversion, has its stated positive lower bound. For arbitrary
encoded real inputs the same construction is an
exact real-arithmetic algorithm; no bit-complexity assertion follows
merely from that convention.

\subsection{Global policy and economic residual certificates}
Set $g_0=0$ and $g_{k+1}=\bar F_Mg_k$, and return $\hat g=g_K$.
Let $g^*$ be the fixed point of the verified operator, implementing
the declared selected competitive equilibrium. The contraction bound
gives
\[
 \|\hat g-F\hat g\|_\infty\leq \bar\rho^K+\delta_M,\qquad
 e:=\|\hat g-g^*\|_\infty
 \leq\frac{\bar\rho^K+\delta_M}{1-\bar\rho}.
 \tag{S12}
\]
For the first inequality,
$\|g_K-\bar F_Mg_K\|\leq
\bar\rho^K\|g_0-\bar F_Mg_0\|\leq \bar\rho^K$, then use
\textup{(S11)}. The second inequality follows by subtracting the two
fixed-point equations and using the exact contraction.

At every state return the finite quantities
\[
 \hat K=\int\hat g(a,e,z,Q_M\mu)\,d(Q_M\mu),\quad
 \hat Q=0.8+\kappa\hat K,\quad
 \hat\pi=\tfrac12\beta\kappa\hat K^2,\quad
 \hat\Phi=Q_M\Phi_{\hat g}(z,Q_M\mu).
 \tag{S13}
\]
The output policy satisfies $\hat g(a,e,z,Q_M\mu)=\hat g(a,e,z,\mu)$
by idempotence of $Q_M$. Therefore the capital-market, price, rebate,
and distribution consistency errors are bounded respectively by
\[
 \frac1{2M},\qquad \frac{\kappa}{2M},\qquad
 \frac{\beta\kappa}{2M},\qquad \frac3{2M}.
 \tag{S14}
\]
The last follows by projecting the future law once and changing its
input law by \textup{(S3)}, using $\operatorname{Lip}_\mu\Phi_{\hat g}\leq2$.
Budget identities hold exactly when consumption is defined from
\textup{(S13)}. It lies in the same strictly positive consumption
interval as in the economic construction. Saving constraints hold
exactly. With the firm's reported physical investment $\hat K$,
the aggregate resource discrepancy is
\[
 \left|\int\hat c\,d\mu+\beta C(\hat K)
              -\int(y_e(z)+a)\,d\mu\right|
 =\beta\hat Q\,|\hat K-K_{\hat g}(\mu)|
 \leq\frac{\beta Q_{\max}}{2M}.
 \tag{S14a}
\]
Here $\hat\pi=\beta\hat Q\hat K-\beta C(\hat K)$ for the stated
quadratic technology. The resource calculation uses that same
reported firm quantity; it does not substitute household saving
into the production cost while using a different firm quantity.

Let $\hat D(b)=\bar E_{\hat g}(b)-
\hat Q(5-y-a-\hat\pi)-(1+\beta\hat Q^2)b$ denote the Euler numerator
using the output functions. A normalized complementarity residual is
$|\hat g-[\hat g+\hat D(\hat g)]_{[0,1]}|$.
This is zero precisely for the correct Euler equality in the interior,
the correct weak inequality at saving zero, and the opposite weak
inequality at saving one. Satisfied slack inequalities contribute zero.
The root of \textup{(S6)} is exactly $\bar F_M\hat g$.
The scalar clipping inequalities thus give
\[
 \sup_{\mathcal D}
 |\hat g-[\hat g+\hat D(\hat g)]_{[0,1]}|
 \leq(1+\beta Q_{\max}^2)(1+m_0)\bar\rho^K
 <1.02\bar\rho^K.
 \tag{S15}
\]
Indeed $[n/d]_{[0,1]}=[b+D(b)/d]_{[0,1]}$, and multiplying a scalar
projected step by $d\geq1$ increases its distance from $b\in[0,1]$
by at most a factor $d$. This check includes both binding constraints.

\subsection{A finite value program and Bellman inequalities}
Use the output policy, prices and distribution transition to compute
$V_H$ by $H$ periods of expected discounted utility and zero terminal
value. Each expectation is the finite sum over $E$ and $Z$.
All intermediate population states are grid supported. This defines
a finite recursive program on the whole domain.

Here are bounds sufficient to include values and Bellman optimality
in the certificate. For any $g\in\mathcal S$, its exact-price,
exact-distribution policy value $V_g$ has absolute value below $13$
and Lipschitz constant at most $6$ in the metric $|a-b|+W_1(\mu,\nu)$
for fixed finite $(e,z)$. In fact $|u(c)|\leq12.5$ and $|u'(c)|\leq5$.
The consumption Lipschitz constants are at most $1.009$ in assets
and $0.0081$ in distributions. Under the common-shock coupling,
the next-state metric is at most three times its current value.
The geometric sum is consequently bounded by
$5(1.009)/(1-3\beta)<6$.

For two policies $g,h$, the current consumption difference is at most
$\beta(Q_{\max}+2\kappa)\|g-h\|_\infty$, and their next asset-plus-
distribution distance is at most $2\|g-h\|_\infty$. Subtract their
policy-value recursions, use the Lipschitz bound $6$, and sum the
discounted error:
\[
 \|V_g-V_h\|_\infty
 \leq
 \frac{5\beta(Q_{\max}+2\kappa)+12\beta}{1-\beta}
 \|g-h\|_\infty
 <0.17\|g-h\|_\infty .
 \tag{S16}
\]
Replacing the exact tuple associated with $\hat g$ by
\textup{(S13)} changes current consumption by at most
$2\beta\kappa/M$ and the next population state by at most $3/(2M)$.
The analogous value-recursion subtraction bounds the induced
policy-value error by
$\beta(9+10\kappa)/\{(1-\beta)M\}<0.1/M$.
The terminal truncation contributes at most $13\beta^H$. Hence
\[
 E_V:=\|V_H-V^*\|_\infty
 \leq0.17e+\frac{0.1}{M}+13\beta^H,
 \tag{S17}
\]
where the economic verification theorem identifies $V_{g^*}=V^*$.

For completeness let $\hat B$ be the true one-period maximization
operator at the returned price and aggregate-transition functions,
with $b$ varying over the full interval $[0,1]$. Comparing it with the
equilibrium Bellman operator, uniformly over that entire interval,
gives
\[
 \|V_H-\hat B V_H\|_\infty
 \leq(1+\beta)E_V+
 (10\beta\kappa+6\beta)e+
 \frac{5\beta\kappa+9\beta}{M}
 \leq0.25e+\frac{0.2}{M}+14\beta^H .
 \tag{S18}
\]
The first additional term bounds the current-price utility change;
the second bounds the continuation-value change when
$W_1(\hat\Phi,\Phi^*)\leq e+3/(2M)$.
The policy Bellman equality error is at most $13\beta^H$, by the
finite-horizon recursion. Consequently the positive part of the
maximal profitable one-step deviation is at most
$0.25e+0.2/M+27\beta^H$. This is a bound for every feasible saving
choice, not a grid-only maximization.

\begin{theorem}[A global finite certificate with operation counts]
\label{thm:finite-solver}
For $0<\varepsilon\leq1$, choose the least nonnegative integers
$K,H$ and the least positive integer $M$ satisfying
\[
 M\geq10/\varepsilon,\qquad
 \bar\rho^K\leq\varepsilon/100,\qquad
 (0.01)^H\leq\varepsilon/100.
 \tag{S19}
\]
The piecewise-affine programs \textup{(S11a), (S13)} and $V_H$ have
all the feasibility, market-clearing, distribution,
Euler-complementarity and Bellman residual blocks bounded by
$\varepsilon$ uniformly over the entire state space and every
primitive in the declared economy class. Their policy error is below
$0.013\varepsilon$ and their value error is below $0.15\varepsilon$.

Let $N=2(M+1)$, $B_k=(2L+3)^{k+1}$,
$\ell_M=1+\lceil\log_2(M+1)\rceil$ and
$\ell_k=1+\lceil\log_2(B_k+1)\rceil$. Define
\[
 C_k=100(k+1)L^2(L+1)^k
 \{N(\ell_M+\ell_k)+B_k\ell_k+1\}.
 \tag{S20}
\]
One query of the complete output tuple, including $V_H$, requires at
most
\[
 C_{\rm enc}
 +100(2L)^{H+1}\{C_K+N\ell_K+1\}
 \tag{S21}
\]
real arithmetic operations and comparisons. For fixed finite $E,Z$
this is polynomial in $1/\varepsilon$, with logarithmic factors.
The output program is finite and defines a function at every
distribution state; it does not enumerate a probability simplex.
The exact-rational implementation and the separately certified
finite-precision alternative below have their stated costs.
\end{theorem}
\begin{proof}
By \textup{(S11)--(S12)},
$e\leq0.012\varepsilon/(1-0.021)<0.013\varepsilon$.
Substitute \textup{(S19)} into \textup{(S14)--(S18)}. The distribution
residual is at most $0.15\varepsilon$ and Bellman optimality violation
is below
$0.25(0.013)\varepsilon+0.02\varepsilon+0.27\varepsilon
<0.30\varepsilon$. All other residuals and the value error have
the asserted bounds.

To count one context construction, build the level-$(k-1)$ policy
pair once at its current context, and once at each of the $L$ future
contexts. Computing the current capital needs $N$ piecewise-affine
evaluations costing at most $O(\ell_{k-1})$ each. For rounding its
future asset positions, locating their two adjacent grid endpoints
by binary search costs at most $O(\ell_M)$ per point. The weights
are then barycentric, not hard cell assignments. Computing each
future price and dividend requires its capital from the returned
level-$(k-1)$ future policy pair: these are an additional $LN$
grid evaluations, not a free by-product of constructing its own
lower-level capital. The factor $100L^2$ in \textup{(S20)} counts
all these evaluations. Forming the union
of the $2L$ future knot lists, its weighted affine coefficients,
and the two inverse lists costs
$O(L B_{k-1}\log(2L B_{k-1}+2))$. All integration, rounding,
splitting, inversion and sorting operations are included.
Thus the context cost $c_k$ satisfies
\[
 c_k\leq(L+1)c_{k-1}
 +100L^2\{N(\ell_M+\ell_k)+B_k\ell_k+1\}.
\]
With a constant-cost zero policy, summing this finite recurrence
gives $c_k\leq C_k$. The current price and population transition
use this constructed policy pair and $N$ further grid evaluations.
Each of the $H$ value dates has at most $2L$ children and the same
bounded local work. Its geometric recursion gives \textup{(S21)}.

For an empirical input with $n_\mu$ atoms, encoding \textup{(S2)}
uses at most $C n_\mu(\ell_M+1)+CN$ operations by barycentric
rounding, or at most $Cn_\mu N$ by the continuous hat formula.
A feature-input query reads $N$ features. Add the cost of reading
the $L^2$ shock matrix and the other primitive data once; these
are $C_{\rm enc}$. Since $K,H=O(\log(1/\varepsilon))$ while $L$
is fixed, every exponential in $K,H$ displayed above is a fixed
power of $1/\varepsilon$. This proves the arithmetic cost claim.
\end{proof}

\paragraph{Bit complexity and a fully finite alternative.}
The polynomial bound just proved uses the declared real-arithmetic
cost model. Exact rational knots give a terminating implementation,
but their bit lengths must be charged; no polynomial bit bound for
nested rational composition is inferred from the arithmetic count.

Here is an alternative with an explicit conservative bit bound.
In \textup{(S6)} use exactly $J$ clipped iterations from zero instead
of the exact piecewise-affine inverse. Call the resulting operator
$T_{M,J}$. The proof of \textup{(S10)} applies to every finite inner
iteration, so it still maps $\mathcal S$ into itself and contracts
by $\bar\rho$. It differs from $\bar F_M$ by at most $m_0^J$. Iterating
it $K$ times therefore gives the same certificates after replacing
$\delta_M$ by $0.02/M+m_0^J$ and the Euler bound by
$1.02(\bar\rho^K+m_0^J)$. Choose the least integers satisfying
\[
 M\geq20/\varepsilon,\qquad
 \bar\rho^K,\ m_0^J,\ (0.01)^H\leq\varepsilon/200
 \tag{S22}
\]
leaves strict residual margins below $\varepsilon/2$.

This alternative can evaluate all hat weights using
$(1-M|a-a_j|)_+$, rather than branching on a rounded asset-cell
index. At a level-$k$ query it makes at most
$N(1+L)+2LJ$ calls at level $k-1$, and at most
$A=100LN^2(J+1)$ remaining arithmetic operations. Hence the complete
query, including prices, distribution and $V_H$, uses at most
\[
 T_{\rm op}=C_{\rm enc}+100N A^{K+1}(2L)^{H+1}+10N
 \tag{S23}
\]
operations. For fixed $L$ this is
$\exp\{O(\log^2(1/\varepsilon))\}$.

The program now consists only of arithmetic, absolute values and
clipping gates; before final normalization, all nonconstant denominators are
$1+\beta Q^2\geq1$. All exact intermediate values are bounded by
$U=100(N+L+1)$. On the enlarged box with denominators at least
$1/2$, each gate is Lipschitz with constant at most $C_0=100U$
in the maximum input error. Request every primitive, the query asset
$a$, all empirical-input atom assets and weights, and every measure
feature to accuracy
\[
 \tau=\frac{\varepsilon}{2000N(2C_0)^{T_{\rm op}+1}}
 \tag{S24}
\]
and round every arithmetic gate to that accuracy. Induction over
the finite computation graph bounds every final coordinate error by
$\delta=\varepsilon/(500N)$; the same induction keeps denominators above
$1/2$ and intermediate values within the enlarged box.
Interval arithmetic can therefore return certified enclosures of
the exact program values. Reported saving approximations may be
clipped to $[0,1]$ without increasing their error.
For a numeric population probability vector, clip negative centre
weights to zero and normalize their sum exactly as rational centre
weights, with no subsequent coordinatewise rounding. Before normalization the
sum is at least $1-N\delta>1/2$. The resulting nonnegative unit-mass
vector has $\ell^1$ error at most $4N\delta$ from the exact vector,
including all preceding rounding errors. Since the grid-state metric
has diameter two, its Wasserstein error is at most this same quantity,
$4N\delta\leq\varepsilon/125$. This prevents coordinatewise errors
from accumulating into an uncharged distribution error. The additional
$O(N)$ clipping, summation and division operations are included in
\textup{(S23)}; their normalization denominator exceeds $1/2$.

The program denotes the exact global function whose invariance
and residual bounds were proved; its certified numeric evaluations
do not redefine it as a discontinuous floating-point policy.
The required precision is
$O(T_{\rm op}\log C_0+\log(1/\varepsilon)+b_{\rm in})$ bits,
where $b_{\rm in}$ includes the primitive, query and empirical-input
bit lengths. Schoolbook
integer arithmetic gives a polynomial overhead in this precision
and $T_{\rm op}$. Thus the fully finite alternative remains
quasipolynomial in $1/\varepsilon$ for fixed $L$, with the charged
polynomial dependence on finite input length and empirical atom
count. Add the declared
cost of obtaining the $N$ input features at this precision, or
their polynomial computation cost for empirical input.
If an input-measure presentation does not supply these computable
features with a stated precision cost, the global mathematical
program and certificate still apply, but no numeric algorithm for
that unencoded measure is asserted.

% Concise paper component. Full optional enhancements remain in inference.tex.
\section{Concrete estimation and numerical likelihood error}
\label{sec:inference}
Fix the known $\beta,P$. The unknown primitive, reference state and moments are
\[
 \theta=(\delta_0,\delta_1,\kappa)\in\Theta
 :=[-0.1,0.1]^2\times[0.0001,0.001],                         \tag{I1}
\]
\[
 \mu^{\rm ref}=\tfrac12\delta_{(L,3/4)}+\tfrac12\delta_{(H,3/4)},\qquad
 p(\theta)=Q_\theta(z_0,\mu^{\rm ref}),\qquad
 m(\theta)=(\delta_0,\delta_1,p(\theta)),                     \tag{I2}
\]
where $z_0$ is fixed. All quantities use the selected regular competitive
equilibrium from Section~\ref{sec:economy}, including endogenous prices and transitions.
\begin{lemma}[Uniform identification]\label{lem:price-identification}
At fixed $\delta$, $p$ is continuous and, for $\kappa_2>\kappa_1$,
\[
 .495(\kappa_2-\kappa_1)\leq p(\delta,\kappa_2)-p(\delta,\kappa_1)
 \leq1.005(\kappa_2-\kappa_1).                              \tag{I3}
\]
At fixed $\kappa$,
\[
 |p(\delta,\kappa)-p(\widetilde\delta,\kappa)|
 \leq.002\|\delta-\widetilde\delta\|_\infty.                  \tag{I4}
\]
Thus, writing $d=\|\delta-\widetilde\delta\|_\infty$,
$k=|\kappa-\widetilde\kappa|$ and $\gamma=.495$,
\[
 \gamma k\leq|p(\theta)-p(\widetilde\theta)|+.002d,\qquad
 \|\theta-\widetilde\theta\|_\infty
 \leq\frac{1.002}{.495}\|m(\theta)-m(\widetilde\theta)\|_\infty,
 \qquad\frac{1.002}{.495}<3.                                \tag{I5}
\]
\end{lemma}
\begin{proof}
Let $K_i=\int g_{\delta,\kappa_i}\,d\mu^{\rm ref}$ at $z_0$.
Proposition~\ref{prop:actual-binding} gives $K_i\geq1/2$ because the
high-income half saves one. Equation \textup{(EC22)} gives
$|K_2-K_1|\leq5(\kappa_2-\kappa_1)$. Hence
\[
 p_2-p_1=(\kappa_2-\kappa_1)K_2+\kappa_1(K_2-K_1)
 \geq(1/2-.001\cdot5)(\kappa_2-\kappa_1).
\]
The same identity with $K_2\leq1$ proves the upper bound. Equation
\textup{(EC23)} proves \textup{(I4)}; both sensitivities imply continuity.
Compare first at fixed $\delta$, then at fixed $\widetilde\kappa$, for
\textup{(I5)}. The income offsets themselves are the other two moments.
\end{proof}

\subsection{Experiment and a certified estimator}
Observe three groups of $r\geq1$ measurements, with known $\sigma>0$:
\[
 X_{zj}=\delta_z+\sigma\xi_{zj}\ (z=0,1),\qquad
 Y_j=p(\theta)+\sigma\xi_{pj},\qquad 1\leq j\leq r,          \tag{I6}
\]
where all $3r$ noises are independent standard normals. $X_{zj}$ is a
measured low-type income minus its baseline one in state $z$; $Y_j$ is a
measured claim price divided by known $\beta$, so its original-price noise
standard deviation is $\beta\sigma$. Each price measurement resets the
population to $\mu^{\rm ref}$ and aggregate state to $z_0$. This specifies
a product Gaussian experiment with $n=3r$, rather than presuming a constant
population along an uncontrolled time series. Equality of observable laws
implies equality of means, so Lemma~\ref{lem:price-identification} proves identification.

The constructed finite solver supplies
\[
 \sup_{\theta\in\Theta}|\widetilde p(\theta)-p(\theta)|\leq\eta. \tag{I7}
\]
If its uniform policy certificate is $e$, its encoding is $Q_M$, and its
certified normalized-price rounding error is $\eta_{\rm fp}$, one may use
\[
 \eta=.001\left(e+\frac1{2M}\right)+\eta_{\rm fp}.           \tag{I8}
\]
Indeed encoding changes capital by at most $1/(2M)$, by the solver's
idempotence and asset Lipschitz properties; policy error adds at most $e$,
and $Q=.8+\kappa K$. Taking $4\mid M$ fixes $\mu^{\rm ref}$ exactly and
removes this encoding term. The solver attains arbitrarily small requested
$\eta$ by its proved resolution, iteration and precision choices.
Thus numerical moments $(\delta_0,\delta_1,\widetilde p)$ have uniform error at most $\eta$.

Let $\overline X_z=r^{-1}\sum_jX_{zj}$, $\overline Y=r^{-1}\sum_jY_j$, and set
\[
 \widehat\delta_z=\operatorname{clip}_{[-.1,.1]}\overline X_z. \tag{I9}
\]
Define $\kappa^\dagger$ to be the unique inverse of
$p(\widehat\delta,\kappa)=\overline Y$ when the target lies between endpoint
prices, and the corresponding endpoint $.0001$ or $.001$ otherwise.
Equation \textup{(I3)} proves existence and uniqueness with this clipping convention.
For $\eta_{\rm bis}>0$, start $[l,u]=[.0001,.001]$ and allow at most
\[
 B=\max\left\{0,\left\lceil\log_2\frac{.0009}{2\eta_{\rm bis}}\right\rceil\right\}
                                                                    \tag{I10}
\]
updates. At midpoint $b$, query $[\widetilde p(\widehat\delta,b)-\eta,
\widetilde p(\widehat\delta,b)+\eta]$. If its upper endpoint is strictly below
$\overline Y$, set $l=b$; if its lower endpoint is strictly above, set $u=b$;
otherwise return $b$. After $B$ updates return the remaining midpoint.
Each discarded half excludes the clipped exact inverse. At an early return,
$|p(\widehat\delta,b)-\overline Y|\leq2\eta$, so \textup{(I3)} gives
distance at most $2\eta/\gamma$, also for a clipped inverse. At exhaustion the
interval midpoint is within $\eta_{\rm bis}$. Therefore the returned estimator obeys
\[
 |\widehat\kappa-\kappa^\dagger|\leq\max\{\eta_{\rm bis},2\eta/\gamma\}
 \leq\eta_{\rm bis}+2\eta/\gamma.                                    \tag{I11}
\]
This needs at most $B$ finite price queries and assumes no monotonicity of
numerical prices. Use rational midpoints and certified comparisons, or
include their rounding in $\eta,\eta_{\rm bis}$. Ideal Gaussian inputs use the declared
real-arithmetic model. For finite encoded sample means of certified accuracy
$\zeta_{\rm in}$, charge their input precision cost and use the supplied
means with every radius $t_r$ replaced by $s_r=t_r+\zeta_{\rm in}$ below.
No finite bit-cost bound for unencoded exact Gaussian reals is asserted.

\begin{theorem}[Finite-sample confidence]\label{thm:gaussian-inference}
For $0<\alpha<1$ put
\[
 t_r=\sigma\sqrt{\frac2r\log\frac6\alpha},\qquad
 R_{\delta,r}=t_r,\qquad
 R_{\kappa,r}=\eta_{\rm bis}+\frac{1.002t_r+2\eta}{.495}.              \tag{I12}
\]
Uniformly over $\theta\in\Theta$, with probability at least $1-\alpha$,
\[
 \|\widehat\delta-\delta\|_\infty\leq R_{\delta,r},\qquad
 |\widehat\kappa-\kappa|\leq R_{\kappa,r}.                   \tag{I13}
\]
Intersect this rectangle with $\Theta$.
\end{theorem}
\begin{proof}
Completing the square gives $\E e^{tG}=e^{t^2/2}$ for a standard normal $G$.
Markov's inequality with $t=x$ and its application to $-G$ give
$\Pr(|G|>x)\leq2e^{-x^2/2}$. Each sample-mean error is
$\sigma G/\sqrt r$. A union bound gives, with probability at least $1-\alpha$,
\[
 \max\{|\overline X_0-\delta_0|,|\overline X_1-\delta_1|,
                         |\overline Y-p(\theta)|\}\leq t_r. \tag{I16}
\]
Clipping cannot increase distance to the true income offsets. The ordered-price
inequality, also at an inverse clipping boundary, gives
$\gamma|\kappa^\dagger-\kappa|\leq
|\overline Y-p(\widehat\delta,\kappa)|\leq1.002t_r$.
Combine with \textup{(I11)}. Encoded means add at most $\zeta_{\rm in}$
by the triangle inequality, proving the stated finite-input version.
\end{proof}
The terms $\eta_{\rm bis}+2\eta/\gamma$ are numerical bias; $t_r$ is Gaussian sampling
error. Choosing $\eta_r,\eta_{{\rm bis},r}=o(r^{-1/2})$, and
$\zeta_{{\rm in},r}=o(r^{-1/2})$ for encoded inputs, makes numerical additions negligible
relative to this sampling radius and gives shrinking uniform confidence bounds.

\subsection{Likelihood error}
The exact observable log likelihood is
\[
 \ell_r(\theta)=-3r\log(\sigma\sqrt{2\pi})-
 \frac1{2\sigma^2}\left\{\sum_{z=0}^1\sum_{j=1}^r(X_{zj}-\delta_z)^2
                         +\sum_{j=1}^r(Y_j-p(\theta))^2\right\}. \tag{I17}
\]
Let $\widetilde\ell_r$ replace only $p$ by $\widetilde p$. Expanding one
Gaussian square with $|\widetilde b-b|\leq\eta$, $|b|\leq B_0$, proves
\[
 |\widetilde\ell(y)-\ell(y)|\leq
 \frac\eta{\sigma^2}(|y|+B_0)+\frac{\eta^2}{2\sigma^2}.       \tag{I18}
\]
Every exact normalized price is at most $B_0=.801$, so deterministically,
for every observed sample and uniformly over the parameter box,
\[
 \sup_{\Theta}|\widetilde\ell_r-\ell_r|\leq\Delta_r
 :=\frac\eta{\sigma^2}\left(\sum_{j=1}^r|Y_j|+rB_0\right)
                       +\frac{r\eta^2}{2\sigma^2}.          \tag{I19}
\]
Since $|Y_j|\leq B_0+\sigma|\xi_{pj}|$, the Gaussian tail and a union bound
over these $r$ noises imply, uniformly over the true primitive, with probability
at least $1-\alpha$,
\[
 \Delta_r/r\leq\frac\eta{\sigma^2}
       \{2B_0+\sigma\sqrt{2\log(2r/\alpha)}\}
                         +\frac{\eta^2}{2\sigma^2}.         \tag{I20}
\]
Thus bounded Gaussian observations have not been assumed. For finite likelihood
comparisons evaluate the quadratic part of \textup{(I17)}; its common normalizer
cancels. Supply known $\sigma$ as positive rational or certified positive interval
input and charge its precision cost, including evaluation of $1/\sigma^2$.
Let each supplied raw $Z^*_{ij}$ have accuracy
$\zeta_{\rm in}$, use group mean bounds $B=(.1,.1,.801)$ and
$\eta_1=\eta_2=0$, $\eta_3=\eta$. Expansion with the numerical mean fixed gives
an additional data-encoding allowance
\[
 D_{\rm in}=\frac{\zeta_{\rm in}}{\sigma^2}
 \sum_{i,j}(|Z^*_{ij}|+\zeta_{\rm in}+B_i+\eta_i)
                         +\frac{3r\zeta_{\rm in}^2}{2\sigma^2}.
\]
Use $\sum|Y_j|\leq\sum|Y_j^*|+r\zeta_{\rm in}$ in \textup{(I19)}, evaluate
the finite sums and squares by rational interval arithmetic, and add its
requested rounding allowance to $\Delta_r+D_{\rm in}$. Charge input precision
and these arithmetic operations. This approximates the ideal Gaussian
likelihood; it does not assign an exact Gaussian density to rounded data.

\section{Counterfactual distributions and welfare}
\label{sec:counterfactual}
A concrete intervention is announced at date zero and permanently replaces
the exogenous income offsets by
\[
 \delta_z^\tau=\operatorname{clip}_{[-.1,.1]}(\delta_z+\tau_z),\qquad
 \theta^\tau=(\delta^\tau,\kappa),\qquad\tau\in\mathbb R^2.  \tag{I24}
\]
Clipping specifies the intervention's version. Initial $(z,\mu)$, discounting,
shock kernels, technology and ownership shares are retained; households
anticipate the new process from date zero. Prices, choices and population
laws are solved again under the same equilibrium selection. This intervention
always stays in $\Theta$ and its parameter map is nonexpansive.
\begin{proposition}[Actual equilibrium counterfactual stability]
\label{prop:counterfactual-stability}
For two admissible primitives with common known $\beta,P$, put
$d=\|\delta-\widetilde\delta\|_\infty$, $k=|\kappa-\widetilde\kappa|$ and
$D=2d+5k$. Starting from the same population and aggregate state, couple the
same shocks. For every aggregate-shock path and every finite $t\geq0$,
\[
 W_1(\mu_t^\theta,\mu_t^{\widetilde\theta})\leq(2^t-1)D.     \tag{I25}
\]
The household optimal values and equal-weight welfare
$W_\theta(z,\mu)=\int V_\theta(a,e,z,\mu)\,d\mu(e,a)$ obey
\[
 \|V_\theta-V_{\widetilde\theta}\|_\infty\leq16d+27k,\qquad
 \sup_{z,\mu}|W_\theta-W_{\widetilde\theta}|\leq16d+27k.      \tag{I26}
\]
\end{proposition}
\begin{proof}
Equations \textup{(EC22)--(EC23)} give $\|g_\theta-g_{\widetilde\theta}\|\leq D$.
If $r_t$ is population Wasserstein distance and $b_t$ individual asset
distance from the same initial household state, the Lipschitz bounds give
\[
 r_{t+1}\leq2r_t+D,\qquad b_{t+1}\leq b_t+r_t+D,\qquad
 r_0=b_0=0.                                                 \tag{I27}
\]
Both are at most $(2^t-1)D$. Capital differs by at most $D+2r_t$, so
with $\kappa_{\max}=.001$,
$|Q_t-\widetilde Q_t|\leq k+\kappa_{\max}(D+2r_t)$ and
$|\pi_t-\widetilde\pi_t|\leq\beta\{k/2+\kappa_{\max}(D+2r_t)\}$.
The budgets, $b_{t+1}\leq b_t+r_t+D$ and $Q_{\max}=.801$ imply
\[
 |c_t-\widetilde c_t|\leq d+\tfrac32\beta k
 +(1+\beta Q_{\max})b_t+\beta(Q_{\max}+4\kappa_{\max})r_t
                         +\beta(Q_{\max}+2\kappa_{\max})D. \tag{I28}
\]
On the consumption interval \textup{(EC9)}, $|u'|\leq5$. Also
$S_\beta:=\sum_{t\geq0}\beta^t(2^t-1)
=\beta/\{(1-2\beta)(1-\beta)\}<.011$,
$1+2\beta Q_{\max}+4\beta\kappa_{\max}<1.02$ and
$\beta(Q_{\max}+2\kappa_{\max})/(1-\beta)<.009$.
Summing \textup{(I28)} with discount $\beta^t$ therefore bounds value differences by
\[
 5\left[\frac{d+\tfrac32\beta k}{1-\beta}
 +\left\{(1+2\beta Q_{\max}+4\beta\kappa_{\max})S_\beta
       +\frac{\beta(Q_{\max}+2\kappa_{\max})}{1-\beta}\right\}D\right]
 \leq6d+k\leq16d+27k.                                      \tag{I29}
\]
Each path follows its own economy's optimal policy, so these are actual
optimal values. Integrating over the common initial population proves welfare.
The resource and dividend accounting remains \textup{(EC3)--(EC6)}.
\end{proof}

Use the confidence radii above and
$\widehat\theta^\tau=(\operatorname{clip}(\widehat\delta+\tau),\widehat\kappa)$.
The finite solver supplies uniform certificates
\[
 \sup_{z,\mu}W_1(\widehat\Phi_\vartheta,\Phi_\vartheta)\leq\nu_\Phi,
 \qquad\|\widehat V_\vartheta-V_\vartheta\|_\infty\leq E_V
 \qquad(\vartheta\in\Theta).                                \tag{I30}
\]
For example $\nu_\Phi=e+3/(2M)$ before rounding, and $E_V$ is the proved
policy, encoding and discounted-tail certificate. Iterate
$\widehat\Phi_{\widehat\theta^\tau}$ from the declared initial population
with common aggregate shocks, giving $\widehat\mu_t^\tau$. Exact transitions
are $2$-Lipschitz, so numerical iteration error satisfies
$v_{t+1}\leq2v_t+\nu_\Phi$, $v_0=0$. On the parameter confidence event,
\textup{(I25)} and nonexpansiveness of the intervention imply
\[
 W_1(\mu_t^{\theta^\tau},\widehat\mu_t^\tau)
 \leq(2^t-1)\{2R_{\delta,r}+5R_{\kappa,r}+\nu_\Phi\}.       \tag{I31}
\]
One may cap the radius at the state-space diameter two. A declared
$L_\psi$-Lipschitz population outcome has mean error at most $L_\psi$ times
this radius; mean mature assets use $L_\psi=1$.

Compute welfare by the finite feature sum
\[
 \widehat W_\vartheta(z,\mu)=
 \int\widehat V_\vartheta(a,e,z,\mu)\,d(Q_M\mu)(e,a),\qquad
 E_W=E_V+3/M.                                               \tag{I32}
\]
The exact value is $6$-Lipschitz in individual initial assets: with population
path fixed, an initial asset difference $b$ stays at most $b$, giving utility
difference at most $5(1+\beta Q_{\max})b/(1-\beta)$, with coefficient below six. The explicit
$Q_M$ coupling preserves type and moves only assets, with mean displacement
at most $1/(2M)$. Hence \textup{(I30)} proves
$|\widehat W_\vartheta-W_\vartheta|\leq E_W$ without an arbitrary-measure
welfare integration oracle. Equal-weight counterfactual welfare and its policy
effect have respective confidence intervals
\[
 \widehat W_{\widehat\theta^\tau}(z,\mu)
 \ \pm\ \{16R_{\delta,r}+27R_{\kappa,r}+E_W\},               \tag{I33}
\]
\[
 \widehat W_{\widehat\theta^\tau}(z,\mu)-\widehat W_{\widehat\theta}(z,\mu)
 \ \pm\ \{32R_{\delta,r}+54R_{\kappa,r}+2E_W\}.              \tag{I34}
\]
Individual value intervals replace $E_W$ by $E_V$. All deterministic bounds
are uniform: the single parameter event gives joint coverage over every initial
state, intervention, aggregate-shock path and finite horizon. The intervals
separate sampling error, price-inversion bias, search resolution, distribution
compression and value truncation for the fully re-solved equilibrium.


\section{Reproducible benchmarks: local and global accuracy}
\label{sec:benchmarks}
The exact-rational programs in \texttt{code/} are
\texttt{benchmark35.py} and \texttt{benchmark35\_value.py}.
Raw machine-readable outputs,
code hashes and resource logs accompany the manuscript. They implement
the finite piecewise-affine construction using Python 3.12.3 exact fractions.
Arithmetic and comparisons are counted; actual numerator and denominator
bit lengths are recorded. Measured times also include Python bookkeeping,
rational gcd work, distribution encoding and preprocessing.

All runs used the same Linux host (Intel Xeon Platinum 8168, 2.70GHz),
one process limited to one CPU and 2GiB RAM, no swap, an internal
55-second combined budget and a 60-second service limit. The common primitives were
\[
 \beta=1/100,\quad\kappa=1/1000,\quad
 (\delta_0,\delta_1)=(-1/10,1/10),\quad
 P=\begin{pmatrix}3/4&1/4\\1/4&3/4\end{pmatrix}.
\]
The policy comparison used $K=2$ updates. Its five initial populations
were the reference half-type population at assets $3/4$; half types at
assets $1/2$; half high types at $1/2$ and low types split equally between
zero and one; all low types at zero; and all high types at one.
Both aggregate states and individual types were evaluated. The seven
asset probes were $0,1/4,1/2,3/4,1$ and $3/4\pm1/64$.

The local comparison builds its own same-resolution reference policy at
$\mu^{\rm ref}$, takes the central own-asset slope at $3/4$ with step $1/64$,
and freezes population dependence. It uses an independent cache and charges
all preprocessing. Both its raw linear extension and clipped version are
reported. The local panel restricts queries to the reference population and
$3/4\pm1/64,3/4$; the wider panel has 140 queries. The global method's
full-domain certificate comes from the proof, rather than these finite panels.

\begin{center}\small
\begin{tabular}{r|rr|rr|r}
$M$ & Global s & Local s & Policy bound & Distribution bound & Clipped local lower\\\hline
16 & 0.0650 & 0.0189 & 0.0017273 & 0.093750 & 0.0014940 \\
32 & 0.0787 & 0.0302 & 0.0010889 & 0.046875 & 0.0021324 \\
64 & 0.1094 & 0.0278 & 0.0007697 & 0.023438 & 0.0024516 \\
\end{tabular}
\end{center}
Times are single-run wall seconds including each method's preprocessing.
Certified upper and lower bounds in the tables are rounded outward.
The policy and distribution columns are rigorous uniform bounds from
\textup{(S12),(S14)}. The last column bounds the clipped local method's error
from below on the finite wider panel: subtract the proved global policy bound
from its exact-rational difference from the computed iterate. It is not an
exhaustive global error bound for the local method. The raw linear extension
violated the asset bound by approximately $0.228916838$; clipping removed
that feasibility violation. Its local-panel comparison and all exact
differences and two-sided bounds are in the data. No simulated-path success
is used as a global certificate.

The complete-tuple experiment also computes capital, both price units,
dividends, the rounded next population and actual $H=2$ policy values,
for both types and aggregate states at the same reference population and
individual asset $3/4$. The clipped local method recomputes its own induced
prices and transition at every continuation. Independent preprocessing
and all value recursion are included. The global value bound is \textup{(S17)}.
\begin{center}\small
\begin{tabular}{r|rr|r}
$M$ & Global tuple s & Local tuple s & Global value bound\\\hline
16 & 0.0528 & 0.0233 & 0.0078437 \\
64 & 0.0839 & 0.0425 & 0.0029934 \\
\end{tabular}
\end{center}
The policy-only runs reached 157/176-bit maximum numerator/denominator
lengths; complete-tuple runs reached 3466/3462 bits. These observed bit lengths
are charged in measured exact-arithmetic times; they are not a polynomial-bit
theorem. A separate $M=8,K=1$ exact first-iterate smoke check and all five
reported runs completed inside their limits. There are no retries or timeouts
in this benchmark set. The four value-query comparisons have the same
finite-query scope, with the uniform value bound enclosing their errors.

The output programs and Sections~\ref{sec:finite-solver}--\ref{sec:counterfactual}
give uniform accuracy choices and separately visible numerical and sampling
errors. The policy error alone is not called the full economic residual:
distribution, market, complementarity and value blocks have their own bounds.

\begin{thebibliography}{9}
\bibitem{Aiyagari} S. Rao Aiyagari (1994).
\emph{Uninsured Idiosyncratic Risk and Aggregate Saving}.
Quarterly Journal of Economics 109(3), 659--684.
\url{https://doi.org/10.2307/2118417}. Background: abstract, p.~659.
\bibitem{KS} Per Krusell and Anthony A. Smith, Jr. (1998).
\emph{Income and Wealth Heterogeneity in the Macroeconomy}.
Journal of Political Economy 106(5), 867--896.
\url{https://doi.org/10.1086/250034}. Background: publisher abstract.
\bibitem{HANK} Greg Kaplan, Benjamin Moll and Giovanni L. Violante (2018).
\emph{Monetary Policy According to HANK}.
American Economic Review 108(3), 697--743.
\url{https://doi.org/10.1257/aer.20160042}. Background: pp.~697--698.
\bibitem{ABRS} Adrien Auclert, Bence Bard\'oczy, Matthew Rognlie and Ludwig Straub (2021).
\emph{Using the Sequence-Space Jacobian to Solve and Estimate Heterogeneous-Agent Models}.
Econometrica 89(5), 2375--2408. \url{https://doi.org/10.3982/ECTA17434}.
Inspected March 2021 author manuscript: abstract, introduction and Sections~4.2,6.
\end{thebibliography}
\end{document}
